\section{Operators on the fibers of a finite product}
\label{sec:perturbations_fibers}

All spaces in this section are finite coordinate spaces $\C^I$ with their
Euclidean inner products, and every coordinate set may be empty unless stated
otherwise.

\begin{definition}[Operators applied on fibers]
    \label{def:euclidean_fiberwise_operator}
    \lean{ContinuousLinearMap.rightFiber, ContinuousLinearMap.rightFiber_apply,
        ContinuousLinearMap.rightFiberwiseMap,
        ContinuousLinearMap.rightFiberwiseMap_apply_apply,
        ContinuousLinearMap.rightFiber_rightFiberwiseMap,
        ContinuousLinearMap.singleRightFiber,
        ContinuousLinearMap.rightFiber_singleRightFiber}
    \leanok
    Let $I,J,S$ be finite sets. For $x\in\C^{I\times S}$ and $s\in S$, the
    fiber $x_s\in\C^I$ is $x_s(i)=x(i,s)$. For a linear operator
    $G:\C^I\to\C^J$, the operator $G\otimes\Id_S:\C^{I\times S}\to\C^{J\times S}$
    applies $G$ to every fiber: $\bigl((G\otimes\Id_S)x\bigr)_s=Gx_s$. For
    $u\in\C^I$ and $s\in S$, the vector $\iota_su\in\C^{I\times S}$ has fiber
    $u$ at $s$ and zero fibers elsewhere.
\end{definition}

\begin{lemma}[Algebra and norms of fiberwise operators]
    \label{thm:euclidean_fiberwise_operator}
    \lean{ContinuousLinearMap.rightFiberwiseMap_add,
        ContinuousLinearMap.rightFiberwiseMap_sub,
        ContinuousLinearMap.rightFiberwiseMap_zero,
        ContinuousLinearMap.rightFiberwiseMap_smul,
        ContinuousLinearMap.rightFiberwiseMap_sum,
        ContinuousLinearMap.rightFiberwiseMap_comp,
        ContinuousLinearMap.rightFiberwiseMap_id,
        ContinuousLinearMap.rightFiberwiseMap_adjoint,
        ContinuousLinearMap.inner_eq_sum_inner_rightFiber,
        ContinuousLinearMap.norm_sq_eq_sum_norm_sq_rightFiber,
        ContinuousLinearMap.inner_rightFiberwiseMap,
        ContinuousLinearMap.norm_singleRightFiber,
        ContinuousLinearMap.rightFiberwiseMap_singleRightFiber,
        ContinuousLinearMap.norm_rightFiberwiseMap_le,
        ContinuousLinearMap.norm_rightFiberwiseMap,
        ContinuousLinearMap.norm_conj_linearIsometryEquiv}
    \leanok
    \uses{def:euclidean_fiberwise_operator}
    The map $G\mapsto G\otimes\Id_S$ is linear, preserves composition,
    identities and adjoints, and
    \begin{align}
        \langle x,y\rangle   & =\sum_{s\in S}\langle x_s,y_s\rangle, &
        \lVert x\rVert^2     & =\sum_{s\in S}\lVert x_s\rVert^2,     &
        \lVert\iota_su\rVert & =\lVert u\rVert.
        \notag
    \end{align}
    Moreover $(G\otimes\Id_S)\iota_su=\iota_sGu$, and
    \begin{align}
        \lVert G\otimes\Id_S\rVert & \leq\lVert G\rVert,
        \label{eq:spectator_fiberwise_norm_le}
    \end{align}
    with equality whenever $S$ is nonempty. Conjugation by a surjective linear
    isometry preserves the operator norm.
\end{lemma}
\begin{proof}
    \leanok
    Linearity, composition and identities are checked on each fiber. Splitting
    the coordinate sums over $I\times S$ by $s$ gives the inner-product and
    norm identities; the adjoint formula follows from them by applying the
    adjoint identity of $G$ in each fiber. Then
    \begin{align}
        \lVert(G\otimes\Id_S)x\rVert^2
         & =\sum_{s\in S}\lVert Gx_s\rVert^2
        \leq\lVert G\rVert^2\sum_{s\in S}\lVert x_s\rVert^2,
    \end{align}
    and for nonempty $S$ the vectors $\iota_su$ give the reverse inequality.
\end{proof}

\begin{lemma}[Operator order after adjoining spectator coordinates]
    \label{thm:ph_spectator_operator_order}
    \lean{ContinuousLinearMap.isSymmetric_rightFiberwiseMap,
        ContinuousLinearMap.isSymmetricProjection_rightFiberwiseMap,
        ContinuousLinearMap.isPositive_rightFiberwiseMap,
        ContinuousLinearMap.rightFiberwiseMap_mono,
        LinearIsometryEquiv.conj_le_conj_iff}
    \leanok
    \uses{thm:euclidean_fiberwise_operator}
    If $G$ is self-adjoint, an orthogonal projection, or positive, so is
    $G\otimes\Id_S$, and $G\le H$ implies $G\otimes\Id_S\le H\otimes\Id_S$. If
    $V$ is a unitary identification of two spaces, then
    $VGV^{-1}\le VHV^{-1}$ if and only if $G\le H$.
\end{lemma}
\begin{proof}
    \leanok
    By the inner-product identity of
    Lemma~\ref{thm:euclidean_fiberwise_operator},
    $\langle(G\otimes\Id_S)x,y\rangle=\sum_s\langle Gx_s,y_s\rangle$, so
    symmetry and positivity hold term by term, and idempotence holds on each
    fiber. Applying positivity to $H-G$ gives the order assertion. Unitary
    conjugation preserves and reflects positivity, hence the order.
\end{proof}

\begin{lemma}[Kernels and ranges of fiberwise operators]
    \label{lem:right-spectator-kernel-projection}
    \lean{ContinuousLinearMap.mem_ker_rightFiberwiseMap_iff,
        ContinuousLinearMap.ker_starProjection_rightFiberwiseMap,
        ContinuousLinearMap.mem_orthogonal_ker_rightFiberwiseMap_iff,
        ContinuousLinearMap.mem_range_rightFiberwiseMap_iff,
        ContinuousLinearMap.mem_orthogonal_range_rightFiberwiseMap_iff,
        ContinuousLinearMap.rightFiberwiseMap_range_starProjection}
    \leanok
    \uses{thm:ph_spectator_operator_order}
    A vector $x$ lies in $\ker(G\otimes\Id_S)$, in its orthogonal complement,
    in $\operatorname{ran}(G\otimes\Id_S)$, or in its orthogonal complement,
    exactly when every fiber $x_s$ lies in the corresponding subspace for $G$.
    Moreover
    \begin{align}
        P_{\ker(G\otimes\Id_S)}               & =P_{\ker G}\otimes\Id_S,              &
        P_{\operatorname{ran}(G\otimes\Id_S)} & =P_{\operatorname{ran}G}\otimes\Id_S.
        \notag
    \end{align}
\end{lemma}
\begin{proof}
    \leanok
    The kernel characterization follows by evaluating each coordinate, and
    the range characterization by choosing a preimage in each fiber; the
    adjoint acts fiberwise, which gives the orthogonal-range criterion. The
    operator $P\otimes\Id_S$ is an orthogonal projection by
    Lemma~\ref{thm:ph_spectator_operator_order} whose range consists of the
    vectors with fibers in the range of $P$; orthogonal projections with the
    same range are equal. The orthogonal complement of the kernel is the
    kernel of $P_{\ker G}^\perp\otimes\Id_S$.
\end{proof}

\begin{lemma}[Joint ranges of fiberwise operators]
    \label{thm:euclidean_fiberwise_projection}
    \lean{ContinuousLinearMap.iSup_range_rightFiberwiseMap_eq,
        ContinuousLinearMap.rightFiberwiseMap_iSup_range_starProjection,
        ContinuousLinearMap.norm_sum_rightFiberwiseMap_range_starProjection_sub_iSup_le}
    \leanok
    \uses{lem:right-spectator-kernel-projection, thm:euclidean_fiberwise_operator}
    For a family $T_j\colon\C^{I_j}\to\C^J$ with varying finite domains and a
    common finite target, the projection onto the joint span satisfies
    \begin{align}
        P_{\sum_j\operatorname{ran}(T_j\otimes\Id_S)}
         & =P_{\sum_j\operatorname{ran}T_j}\otimes\Id_S.
    \end{align}
    The index set may be infinite, with each sum read as the linear span. For a
    finite family,
    \begin{align}
        \left\|\sum_jP_{\operatorname{ran}(T_j\otimes\Id_S)}
        -P_{\sum_j\operatorname{ran}(T_j\otimes\Id_S)}\right\|
        \leq
        \left\|\sum_jP_{\operatorname{ran}T_j}
        -P_{\sum_j\operatorname{ran}T_j}\right\|.
    \end{align}
    These identities give the spectator projection conjugation used in
    \cite[Lemma~commutation(i)]{Nachtergaele1996Spectral}.
\end{lemma}
\begin{proof}\leanok
    For a joint span, orthogonality to every summand and membership of every
    fiber commute as universal quantifiers, and taking orthogonal complements
    again proves the joint-range identity. The estimate follows from linearity
    of $G\mapsto G\otimes\Id_S$ and the norm
    bound~(\ref{eq:spectator_fiberwise_norm_le}).
\end{proof}

\begin{lemma}[Gaps of fiberwise operators]
    \label{thm:perturbations_fiberwise_gap}
    \lean{ContinuousLinearMap.singleRightFiber_mem_orthogonal_ker_rightFiberwiseMap_iff,
        ContinuousLinearMap.norm_gap_rightFiberwiseMap,
        ContinuousLinearMap.norm_gap_rightFiberwiseMap_iff}
    \leanok
    \uses{lem:right-spectator-kernel-projection, thm:euclidean_fiberwise_operator}
    Let $G:\C^I\to\C^J$ and $\delta\ge0$. If $\delta\lVert u\rVert\le\lVert Gu\rVert$
    for every $u\perp\ker G$, then
    $\delta\lVert x\rVert\le\lVert(G\otimes\Id_S)x\rVert$ for every
    $x\perp\ker(G\otimes\Id_S)$. For nonempty $S$ the converse holds as well. No
    positivity or symmetry of $G$ is assumed.
\end{lemma}
\begin{proof}
    \leanok
    For $x\perp\ker(G\otimes\Id_S)$ every fiber satisfies $x_s\perp\ker G$, and
    summing $\delta^2\lVert x_s\rVert^2\le\lVert Gx_s\rVert^2$ over $s$ gives the
    first claim. Conversely, $\iota_su\perp\ker(G\otimes\Id_S)$ exactly when
    $u\perp\ker G$, with $\lVert\iota_su\rVert=\lVert u\rVert$ and
    $\lVert(G\otimes\Id_S)\iota_su\rVert=\lVert Gu\rVert$.
\end{proof}

\begin{lemma}[Gaps under a fiberwise conjugacy]
    \label{lem:perturbations_fiberwise_conjugate_gap}
    \lean{ContinuousLinearMap.smul_orthogonal_ker_projection_le_of_conj_rightFiberwiseMap,
        ContinuousLinearMap.forall_hasEigenvalue_of_conj_rightFiberwiseMap}
    \leanok
    \uses{lem:right-spectator-kernel-projection, lem:gap_isometric_conjugation,
        lem:gap_operator_form}
    Let $E$ be a finite-dimensional complex Hilbert space, $U:E\to\C^{I\times S}$ a
    surjective linear isometry, $G$ an operator on $E$ and $H$ an operator on $\C^I$ with
    $UGU^{-1}=H\otimes\Id_S$. Let $P_G$ and $P_H$ be the orthogonal projections onto
    $(\ker G)^\perp$ and $(\ker H)^\perp$, and $\kappa\in\R$. If $\kappa P_H\le H$,
    then $\kappa P_G\le G$. Consequently, if $H\ge0$, $\kappa>0$ and every eigenvalue
    of $H$ other than $0$ is at least $\kappa$, then every eigenvalue of $G$ other than
    $0$ is at least $\kappa$.
\end{lemma}
\begin{proof}
    \leanok
    By Lemma~\ref{lem:gap_isometric_conjugation}, $U$ conjugates the projection onto
    $\ker G$ to the projection onto $\ker(H\otimes\Id_S)$, which is the extension of
    the projection onto $\ker H$; hence $UP_GU^{-1}=P_H\otimes\Id_S$. Extending
    $\kappa P_H\le H$ over the fibers gives $\kappa P_H\otimes\Id_S\le H\otimes\Id_S$,
    and conjugating back by $U$ gives $\kappa P_G\le G$. For the second claim,
    Lemma~\ref{lem:gap_operator_form} turns the eigenvalue bound of $H$ into
    $\kappa P_H\le H$, and the bound $\kappa P_G\le G$ back into the eigenvalue
    bound of $G$.
\end{proof}

\begin{lemma}[Quadratic-form bound on fibers]
    \label{lem:right_spectator_quadratic_form_gap}
    \lean{ContinuousLinearMap.quadraticForm_sq_ge_rightFiberwiseMap,
        LinearMap.IsPositive.quadraticForm_sq_ge_rightFiberwiseMap_of_forall_hasEigenvalue}
    \leanok
    \uses{lem:gap_norm_to_quadratic_form, thm:euclidean_fiberwise_operator}
    Let $G$ act on $\C^I$. If $G^2\geq\gamma G$ as a quadratic form, the same
    inequality holds for $G\otimes\Id_S$. In particular, if $G\ge0$ and every
    eigenvalue of $G$ other than $0$ is at least $\gamma\ge0$, this inequality holds
    for $G\otimes\Id_S$.
\end{lemma}
\begin{proof}
    \leanok
    Both quadratic forms split into finite sums over $s$, so the inequality
    for $G$ gives
    \begin{align}
        \gamma\sum_s\operatorname{Re}\langle Gx_s,x_s\rangle
         & \leq \sum_s\operatorname{Re}\langle Gx_s,Gx_s\rangle,
    \end{align}
    and Lemma~\ref{lem:gap_norm_to_quadratic_form} gives the last statement.
\end{proof}

\begin{definition}[Finite dependent direct sums]
    \label{def:mpo_sigma_fiber_operator}
    \lean{ContinuousLinearMap.sigmaFiber, ContinuousLinearMap.sigmaFiberwiseMap,
        ContinuousLinearMap.sigmaFiberwiseMap_apply_apply,
        ContinuousLinearMap.sigmaFiber_sigmaFiberwiseMap,
        ContinuousLinearMap.singleSigmaFiber,
        ContinuousLinearMap.sigmaFiber_singleSigmaFiber}
    \leanok
    Let $Q$ be finite, let $J_q$ be finite for each $q\in Q$, and set
    $E_q=\C^{J_q}$ and $E=\C^{\coprod_{q\in Q}J_q}$. For $v\in E$, its
    restriction $v_q\in E_q$ is $v_q(i)=v(q,i)$. A family of linear maps
    $F_q:E_q\to E_q$ defines $F_\Sigma:E\to E$ by $(F_\Sigma v)_q=F_qv_q$.
    For $q\in Q$ and $u\in E_q$, the vector $\iota_qu\in E$ has restriction
    $u$ at $q$ and zero at every other label.
\end{definition}

\begin{lemma}[Inner products, norms and kernels in a finite direct sum]
    \label{thm:mpo_sigma_fiber_inner_norm}
    \lean{ContinuousLinearMap.inner_eq_sum_inner_sigmaFiber,
        ContinuousLinearMap.norm_sq_eq_sum_norm_sq_sigmaFiber,
        ContinuousLinearMap.norm_singleSigmaFiber,
        ContinuousLinearMap.inner_singleSigmaFiber,
        ContinuousLinearMap.sigmaFiberwiseMap_singleSigmaFiber,
        ContinuousLinearMap.mem_ker_sigmaFiberwiseMap_iff,
        ContinuousLinearMap.mem_orthogonal_ker_sigmaFiberwiseMap_iff}
    \leanok
    \uses{def:mpo_sigma_fiber_operator}
    For $v,w\in E$, $q\in Q$ and $u\in E_q$,
    \begin{align}
        \langle v,w\rangle       & =\sum_{q\in Q}\langle v_q,w_q\rangle,        &
        \|v\|^2                  & =\sum_{q\in Q}\|v_q\|^2,              \notag   \\
        \|\iota_qu\|             & =\|u\|,                                      &
        \langle\iota_qu,v\rangle & =\langle u,v_q\rangle,             \notag
    \end{align}
    and $F_\Sigma\iota_qu=\iota_qF_qu$. Moreover $v\in\ker F_\Sigma$ if and only
    if $v_q\in\ker F_q$ for every $q$, and $v\perp\ker F_\Sigma$ if and only if
    $v_q\perp\ker F_q$ for every $q$.
\end{lemma}
\begin{proof}
    \leanok
    Group the coordinates with the same label. If $v\perp\ker F_\Sigma$, testing
    against $\iota_qu$ for $u\in\ker F_q$ gives $\langle u,v_q\rangle=0$.
    Conversely, for $w\in\ker F_\Sigma$ each summand of
    $\langle w,v\rangle=\sum_q\langle w_q,v_q\rangle$ vanishes.
\end{proof}

\begin{lemma}[Norm gaps in a finite direct sum]
    \label{thm:mpo_sigma_fiber_gap_iff}
    \lean{ContinuousLinearMap.norm_gap_sigmaFiberwiseMap,
        ContinuousLinearMap.norm_gap_sigmaFiberwiseMap_iff}
    \leanok
    \uses{thm:mpo_sigma_fiber_inner_norm}
    Let $\delta\geq0$. Then $\delta\|v\|\leq\|F_\Sigma v\|$ for every
    $v\perp\ker F_\Sigma$ if and only if $\delta\|u\|\leq\|F_qu\|$ for every
    $q\in Q$ and $u\perp\ker F_q$. No nonemptiness, positivity, or symmetry
    assumption is imposed.
\end{lemma}
\begin{proof}
    \leanok
    Summing the inequalities $\delta^2\|v_q\|^2\leq\|F_qv_q\|^2$ gives one
    implication; applying the bound to $\iota_qu$ gives the other.
\end{proof}

\begin{definition}[Operators with dependent spectator coordinates]
    \label{def:mpo_dependent_spectator_operator}
    \lean{ContinuousLinearMap.dependentRightFiber, ContinuousLinearMap.dependentRightFiberwiseMap,
        ContinuousLinearMap.dependentRightFiberwiseMap_apply_apply,
        ContinuousLinearMap.dependentRightFiber_dependentRightFiberwiseMap}
    \leanok
    \uses{def:mpo_sigma_fiber_operator, def:euclidean_fiberwise_operator}
    Let $Q$ be a finite set, and let $I_q$ and $S_q$ be finite sets for each
    $q\in Q$. For $v\in\C^{\coprod_q(I_q\times S_q)}$, $q\in Q$, and $s\in S_q$,
    define $v_{q,s}\in\C^{I_q}$ by $v_{q,s}(i)=v(q,i,s)$. Given linear
    operators $G_q$ on $\C^{I_q}$, the operator $H_S$ is the direct sum of the
    $G_q\otimes\Id_{S_q}$:
    \begin{align}
        (H_Sv)_{q,s} & =G_qv_{q,s}.
        \label{eq:mpo_dependent_spectator_action}
    \end{align}
\end{definition}

\begin{lemma}[Kernels and norms with dependent spectators]
    \label{thm:mpo_dependent_spectator_kernel}
    \lean{ContinuousLinearMap.mem_ker_dependentRightFiberwiseMap_iff,
        ContinuousLinearMap.mem_orthogonal_ker_dependentRightFiberwiseMap_iff,
        ContinuousLinearMap.norm_sq_eq_sum_norm_sq_dependentRightFiber}
    \leanok
    \uses{def:mpo_dependent_spectator_operator, thm:mpo_sigma_fiber_inner_norm,
        lem:right-spectator-kernel-projection}
    A vector $v$ lies in $\ker H_S$, respectively in its orthogonal complement,
    if and only if $v_{q,s}$ lies in $\ker G_q$, respectively in its orthogonal
    complement, for every $q\in Q$ and $s\in S_q$. Moreover
    $\|v\|^2=\sum_{q\in Q}\sum_{s\in S_q}\|v_{q,s}\|^2$.
\end{lemma}
\begin{proof}
    \leanok
    Combine the corresponding statements for finite direct sums and for
    operators applied on fibers.
\end{proof}

\begin{lemma}[Common norm gaps with dependent spectators]
    \label{thm:mpo_dependent_spectator_gap_iff}
    \lean{ContinuousLinearMap.norm_gap_dependentRightFiberwiseMap_iff,
        ContinuousLinearMap.norm_gap_dependentRightFiberwiseMap_iff_dependentRightFiberwiseMap}
    \leanok
    \uses{thm:mpo_sigma_fiber_gap_iff, thm:perturbations_fiberwise_gap}
    Let $\delta\geq0$ and let every $S_q$ be nonempty. Then
    $\delta\|v\|\leq\|H_Sv\|$ for every $v\perp\ker H_S$ if and only if
    $\delta\|u\|\leq\|G_qu\|$ for every $q\in Q$ and $u\perp\ker G_q$. Hence
    for nonempty $S_q$ and $T_q$ the operators $H_S$ and $H_T$ built from the
    same family $(G_q)$ have the same such bounds, whatever the cardinalities of
    $S_q$ and $T_q$.
\end{lemma}
\begin{proof}
    \leanok
    Apply Lemma~\ref{thm:mpo_sigma_fiber_gap_iff} to the family
    $G_q\otimes\Id_{S_q}$ and Lemma~\ref{thm:perturbations_fiberwise_gap} to
    each of its members.
\end{proof}

\begin{lemma}[Dependent orthogonal projections]
    \label{thm:mpo_dependent_orthogonal_projections}
    \lean{ContinuousLinearMap.isSymmetricProjection_sigmaFiberwiseMap,
        ContinuousLinearMap.isSymmetricProjection_dependentRightFiberwiseMap,
        LinearIsometryEquiv.conj_eq_dependentRightFiberwiseMap_of_ker_iff}
    \leanok
    \uses{thm:mpo_dependent_spectator_kernel}
    Let $P_q$ be orthogonal projections on $\C^{I_q}$. Then
    $\bigoplus_q(P_q\otimes\Id_{S_q})$ is an orthogonal projection.
    If $P$ is an orthogonal projection and
    $U:\mathcal H\to\C^{\coprod_q(I_q\times S_q)}$ is unitary with
    $v\in\ker P$ exactly when $(Uv)_{q,s}\in\ker P_q$ for every $q,s$,
    then $UPU^{-1}=\bigoplus_q(P_q\otimes\Id_{S_q})$.
\end{lemma}
\begin{proof}
    \leanok
    The inner product is the sum of the fiber inner products, so symmetry
    and idempotence hold on each fiber. The two conjugated projections
    have the same kernel, hence the same range, the orthogonal complement
    of that kernel. Orthogonal projections with the same range are equal.
\end{proof}

\begin{definition}[Euclidean extension by zero]
    \label{def:physical_euclidean_extension_zero}
    \lean{Function.Embedding.euclideanExtendByZero,
        Function.Embedding.euclideanExtendByZero_apply_image,
        Function.Embedding.euclideanExtendByZero_apply_not_mem_range}
    \leanok
    For an injection $e:I\hookrightarrow J$ of finite sets, let
    $E_e:\mathbb C^I\to\mathbb C^J$ extend each coefficient family by zero:
    \begin{align}
        (E_e\xi)(j) & =
        \begin{cases}
            \xi(i), & j=e(i),       \\
            0,      & j\notin e(I).
        \end{cases}
    \end{align}
\end{definition}

\begin{lemma}[Extension by zero is an isometry]
    \label{lem:physical_euclidean_extension_zero_isometry}
    \lean{Function.Embedding.norm_euclideanExtendByZero}
    \leanok
    \uses{def:physical_euclidean_extension_zero}
    For every injection $e:I\hookrightarrow J$ of finite sets, the map $E_e$
    is linear and $\lVert E_e\xi\rVert=\lVert\xi\rVert$ for every
    $\xi\in\mathbb C^I$.
\end{lemma}
\begin{proof}
    \leanok
    Since $(E_e\xi)(j)=0$ for $j\notin e(I)$ and $e$ is injective,
    \begin{align}
        \sum_{j\in J}\lvert(E_e\xi)(j)\rvert^2
         & =\sum_{i\in I}\lvert(E_e\xi)(e(i))\rvert^2
        =\sum_{i\in I}\lvert\xi(i)\rvert^2.
    \end{align}
\end{proof}

\begin{lemma}[Fiberwise operators and extension by zero]
    \label{lem:physical_spectator_embedding_transport}
    \lean{ContinuousLinearMap.rightFiberwiseMap_extendByZero}
    \leanok
    \uses{def:physical_euclidean_extension_zero, def:euclidean_fiberwise_operator}
    Let $e:I\hookrightarrow J$ and $g:S\hookrightarrow T$ be injections of
    finite sets, and let $G$ and $H$ be linear operators on
    $\mathbb C^I$ and $\mathbb C^J$. Suppose $HE_e=E_eG$. Then
    \begin{align}
        (H\otimes\Id_T)E_{e\times g} & =E_{e\times g}(G\otimes\Id_S).
    \end{align}
\end{lemma}
\begin{proof}
    \leanok
    The fiber at $g(s)$ of an extended vector is the extension of its fiber at
    $s$, and the fibers outside $g(S)$ vanish. Apply the intertwining relation
    to the first fibers and linearity to the zero fibers.
\end{proof}
